\documentclass[11pt]{article}

\usepackage[margin=1.1in]{geometry}
\usepackage[utf8]{inputenc}
\usepackage[T1]{fontenc}
\DeclareUnicodeCharacter{03B8}{\ensuremath{\theta}}
\usepackage{amsmath,amssymb}
\usepackage{booktabs}
\usepackage{listings}
\usepackage{xcolor}
\usepackage{microtype}
\usepackage[hidelinks]{hyperref}

% Same plain Julia listing style as model_definition.tex.
\lstdefinelanguage{Julia}{
  keywords={abstract,type,struct,mutable,function,end,if,else,elseif,for,
            while,begin,return,using,import,export,const,where,in,do,let,
            quote,macro,module,true,false,nothing,missing,local,global},
  sensitive=true,
  morecomment=[l]{\#},
  morestring=[b]",
}
\lstset{
  language=Julia,
  basicstyle=\ttfamily\small,
  keywordstyle=\bfseries,
  commentstyle=\color{black!55}\itshape,
  stringstyle=\color{black!70},
  showstringspaces=false,
  columns=fullflexible,
  keepspaces=true,
  breaklines=true,
  breakatwhitespace=true,
  frame=single,
  framerule=0.3pt,
  xleftmargin=1em,
  xrightmargin=1em,
  aboveskip=0.8\baselineskip,
  belowskip=0.8\baselineskip,
  literate={θ}{{$\theta$}}1
           {∂}{{$\partial$}}1
           {≤}{{$\le$}}1
           {≥}{{$\ge$}}1
           {∈}{{$\in$}}1
           {∞}{{$\infty$}}1
           {·}{{$\cdot$}}1
           {—}{{---}}1
           {–}{{--}}1,
}

\newcommand{\E}{\mathbb{E}}
\newcommand{\code}[1]{\texttt{#1}}

\title{Concourse: The Event Loop}
\author{Concourse research notes}
\date{July 17, 2026}

\begin{document}
\maketitle
\tableofcontents

%==============================================================================
\section{What this document is}

\code{model\_definition.tex}~\S8 queued four middle-layer questions. Design
discussion on July 17, 2026 settled three and framed the fourth:

\begin{description}
\item[D1 (decided).] \code{fire} emits enable/disable/re-enable deltas; the
  full \code{enabled(state)} recompute is retained as a debug-mode
  consistency check, not measured against --- deltas are chosen for
  correctness of the order of operations, not for speed.
\item[D2 (open; proposal in \S\ref{sec:cascade}).] The cascade worklist must
  settle to a fixed point; the policy that makes the fixed point unique is
  proposed here and awaits a decision.
\item[D3 (decided).] The interpreter owns the marked record. CompetingClocks
  samples, and only samples; its watchers observe sampling, and the marked
  record is not a sampling artifact.
\item[D4 (decided).] Sampler choice belongs to CompetingClocks'
  \code{SamplerBuilder}: Concourse describes the clock-key type and the
  distributions in play, and the builder picks. Any sampler valid for those
  distributions must work; none is load-bearing.
\end{description}

Charter claims F9--F11 (\S\ref{sec:charter}) extend the falsification
charter of \code{model\_definition.tex}~\S7.

%==============================================================================
\section{D1: \code{fire} emits clock deltas}
\label{sec:deltas}

\subsection{Why the diff loses}

Beyond ordering, an enabled-set diff is \emph{semantically blind to
re-enablings}. When processor sharing re-evaluates every resident clock on an
arrival, or a preemption changes which distribution a still-enabled key
carries, the enabled set before and after the firing is identical --- the
diff sees nothing, while the sampler must receive \code{reenable!}. Deltas
carry information the diff cannot express. The diff survives only as the
debug oracle for the one thing it can check: membership.

\subsection{The delta vocabulary}

The contract's \code{fire} is $\theta$-free and pure, so deltas cannot carry
distributions --- only keys and operations. The interpreter derives each
distribution from \code{clock\_distribution} at the post-firing state and
each anchor from \code{st.te} (amendment A1):

\begin{lstlisting}
const ClockDelta = Tuple{Symbol,ClockKey}   # (:enable | :disable | :reenable, key)
# fire returns (new_state, deltas::Vector{ClockDelta}); emission order is
# authoritative and is applied to the sampler verbatim.
\end{lstlisting}

This is ClockGradients' existing \code{fire\_changes(model, state, key)}
seam with the changes type made concrete (and the CD-1 time argument
applied to it as well); \code{enabled\_update} consumes the same deltas, so
the estimators' incremental path and the interpreter share one vocabulary.
The fired clock itself is implicit: the interpreter calls
\code{fire!(ctx, key, t)} before applying deltas, so \code{fire} never emits
a delta for its own key.

\subsection{The debug oracle}

After every firing, debug mode asserts two things:

\begin{enumerate}
\item \emph{Membership:} \code{enabled(m, st')} equals the previous enabled
  set plus emitted \code{:enable}s minus emitted \code{:disable}s (and the
  fired key). This is the complete-recompute check, and it is the
  specification --- \code{enabled} remains the meaning, deltas the
  implementation.
\item \emph{Anchors:} the sampler's per-clock enabling times equal
  \code{st.te} --- the A1 invariant. This is what checks re-enablings, which
  membership cannot see.
\end{enumerate}

\subsection{Refinement forced by F2: deltas are derived, not emitted}
\label{sec:derived}

Running charter claim F2 (add reneging touching only the family registry and
the surface constructor) falsified D1 \emph{as first implemented}: patience
clocks are born when a job enters a buffer and die when it leaves, and with
deltas hand-emitted inside \code{settle!} and the family \code{fire}
functions, adding the family would have meant threading emission code
through every state-transition site — interpreter surgery, which is exactly
what F2 forbids. The repair keeps D1's decision and changes its mechanism:

\begin{itemize}
\item \emph{Membership deltas are derived, not emitted.} The cascade already
  knows which stations it touched. After the cascade, for each touched
  station and each family, the family's per-station enabled keys are diffed
  between the old and new states (both exist; \code{fire} is pure). A family
  is now exactly four methods — \code{family\_station\_keys!},
  \code{family\_law}, \code{family\_fire!}, \code{family\_memory} — and
  clock lifecycle follows from membership alone.
\item \emph{Locality is preserved.} The user's original argument for deltas
  over a global diff was cost on large networks; the derived diff runs only
  over the touched stations' clocks, so the asymptotics stand. Explicit
  \code{:reenable} deltas remain available for the processor-sharing path
  (phase 2), which membership cannot see.
\item \emph{The A1 bookkeeping collapses into one place.} Enabling times and
  banked ages are updated in the same derivation pass, driven by the diffs
  and the per-family memory policy, instead of at every site that moves a
  job.
\item \emph{A correctness dividend.} A clock enabled and disabled within one
  firing (a job dispatched and immediately preempted in the same cascade)
  now nets out of the diff and never reaches the sampler — the GSMP has no
  zero-duration clocks, which the emission order previously only
  approximated by a disable/enable pair.
\end{itemize}

This is the falsification charter working as designed: the claim failed
against the first implementation, and what survived is a better mechanism
under the same decision. After the repair, the \code{:patience} addition
itself changed only the registry and the surface fields, and M/M/1+M
matches the Erlang-A birth--death form with the membership oracle running
through every reneging cascade.

%==============================================================================
\section{D2 proposal: the cascade worklist}
\label{sec:cascade}

\subsection{The problem}

Within one firing at time $t$, instantaneous consequences chain: a deposit
triggers \code{settle!} at the destination; a departure from a full buffer
lets a blocked upstream server transfer in, which frees that server, which
admits from its buffer, which may unblock stations further upstream. The
cascade must reach a fixed point --- no station can dispatch, preempt, or
unblock --- before the sampler sees any delta. The requirement is that the
fixed point be \emph{unique and deterministic}, because \code{fire} is the
replayed, differentiated object.

The GSPN literature has this exact problem as immediate transitions and
vanishing markings, resolved there by priorities and weights (Zimmermann
Ch.~5--6). The proposal below is that resolution, restated for queues.

\subsection{Split contention from propagation}

The order of cascade processing matters in exactly one place: when two
consumers compete for one freed resource --- two blocked upstream servers and
one freed buffer slot; two idle servers and one arriving job is already
settled by the discipline's \code{select}. Everywhere else the steps merely
propagate.

\begin{description}
\item[Contention is model semantics.] Who gets a freed slot is an observable
  modeling choice, not an implementation detail, so it must be a declared
  policy value on the station, parallel to the discipline:
\begin{lstlisting}
UnblockPolicy:
    FCFSUnblock()          # longest-blocked first (default)
    PriorityUnblock(by)    # by a ScalarExpr over the blocked job's marks
    RandomUnblock()        # a recorded choice, via the A4 draw source
\end{lstlisting}
  \code{FCFSUnblock} needs the blocking order, so the blocked-server set in
  \code{QueueState} is an ordered vector, not a set. \code{RandomUnblock}
  consumes a recorded draw and so stays inside the A4 randomness census.
\item[Propagation must be order-free.] With contention delegated to policy,
  the claim is that the remaining cascade relation is \emph{confluent}: any
  fair processing order reaches the same fixed point. Termination comes from
  check C3 (acyclic blocking chains): each unblocking transfer moves a job
  strictly downstream in the blocking DAG, so the cascade measure is
  well-founded. Confluence plus termination give a unique normal form
  (Newman's lemma); the worklist order is then an implementation choice, not
  a semantic one.
\end{description}

The canonical order, so that the implementation is deterministic even if the
confluence claim fails: a FIFO worklist of stations to settle, seeded with
the destination station then the origin station of the firing --- the
one-hop order \code{queue\_layers.tex}~\S3.2 fixed by fiat, generalized. A
debug-build fuel counter bounds the cascade at (jobs + servers) steps and
turns nontermination into an error rather than a hang.

\subsection{What falsifies it}

F9 in the charter: run the same trajectories with FIFO and LIFO worklists;
\code{states\_equal} at every index or the confluence claim is dead, and the
worklist order gets promoted into the IR's semantics the way the one-hop
order already is. Either outcome is a result; the test is cheap.

%==============================================================================
\section{D3: the interpreter owns the record}

The marked record $(k_i, t_i, a_i)$ is the primary observable of the whole
library, and its auxiliary-draw column exists because the \emph{interpreter}
supplies the draw source (A4) --- the sampler never sees those draws, so a
sampler-side watcher could never capture them. The interpreter appends one
entry per firing: the key, the time, and the draws consumed by that firing
in consumption order. CompetingClocks' watchers remain what they are ---
instrumentation of sampling --- and Concourse attaches none of the record
path to them. \code{gradient\_record} gets one method, from the marked
record type, per \code{model\_definition.tex}~\S6.

One replay detail follows from record ownership: during estimation,
\code{fire} must read firing $i$'s recorded draws, and the contract does not
pass $i$. \code{QueueState} therefore carries a firing counter
(\code{nfired::Int}, incremented by every \code{fire}), and the replay
binding indexes the record with it. A pure counter in a pure state; the
same device as \code{next\_id}.

%==============================================================================
\section{D4: the builder chooses the sampler}

Concourse tells \code{SamplerBuilder} the clock-key type and lets
CompetingClocks choose the sampler for the distributions in play; the
default is whatever the builder decides, and Concourse hard-codes no method.
The design consequence is a claim, not a configuration: \emph{every} sampler
valid for the model's distribution set must produce statistically identical
results, because the interpreter speaks only the blessed
\code{SamplingContext} surface (\code{enable!}, \code{disable!},
\code{reenable!}, \code{next}, \code{fire!}) and no model code reaches
around it. Estimator-driven overrides (the carry coupling for IPA) are
passed through as builder arguments by the estimation entry points, not
chosen by the model.

%==============================================================================
\section{The interpreter, whole}

With D1--D4 fixed, the middle layer is small enough to state completely.
Julia-flavored pseudocode; error handling and the debug oracle elided:

\begin{lstlisting}
function simulate(m::QueueGSMP, θ, horizon, seed)
    st  = initial_state(m)
    ctx = SamplingContext(build(m), seed)     # D4: builder chooses
    rec = MarkedRecord(clockkeytype(m))
    for k in enabled(m, st)                   # initial enablings
        enable!(ctx, k, clock_distribution(m, θ, k, st), st.te[k], 0.0)
    end
    while true
        (t, key) = next(ctx, st.time)
        (key === nothing || t > horizon) && break
        fire!(ctx, key, t)                    # realize the winner's draw
        draws = draw_source(ctx, key)         # A4: keyed, recording
        st, deltas = fire_changes(m, st, key, t, draws)   # cascade inside
        push!(rec, key, t, consumed(draws))   # D3: interpreter-owned
        for (op, k) in deltas                 # D1: emission order verbatim
            op === :disable  && disable!(ctx, k, t)
            op === :enable   && enable!(ctx, k, clock_distribution(m, θ, k, st), st.te[k], t)
            op === :reenable && reenable!(ctx, k, clock_distribution(m, θ, k, st), st.te[k], t)
        end
    end
    close_record!(rec, min(horizon, st.time))
end
\end{lstlisting}

Everything model-specific is inside \code{fire\_changes} (families,
disciplines, kernels, the cascade); everything stochastic is inside
\code{ctx} and \code{draws}; everything observable is \code{rec}. The loop
itself should never grow.

%==============================================================================
\section{Charter extensions}
\label{sec:charter}

Continuing the numbering of \code{model\_definition.tex}~\S7:

\begin{description}
\item[F9 (cascade confluence).] With contention resolved by
  \code{UnblockPolicy}, the cascade fixed point is order-independent.
  \emph{Test:} FIFO vs.\ LIFO worklists over random finite-buffer
  \code{:block} networks; \code{states\_equal} at every firing index.
  Falsified $\Rightarrow$ the worklist order becomes IR semantics.
\item[F10 (sampler agnosticism).] No charter result depends on the sampler.
  \emph{Test:} run the F1--F7 oracle comparisons under at least two builder
  outcomes (e.g.\ first-to-fire and next-reaction); all pass under both.
\item[F11 (delta fidelity).] Emitted deltas agree with the recompute
  oracle. \emph{Test:} debug-mode membership and anchor assertions
  (\S\ref{sec:deltas}) hold over long mixed runs including preemption ---
  and, in phase 2, processor sharing, which is the case only the anchor
  check can see.
\end{description}

\end{document}
